\documentclass[11pt]{article}
\usepackage[T1]{fontenc}
\usepackage[utf8]{inputenc}
\usepackage[margin=2.2cm]{geometry}
\usepackage{booktabs}
\usepackage{graphicx}
\usepackage{float}
\usepackage{longtable}
\usepackage{microtype}
\graphicspath{{../../../outputs/reports/graphpro_dataset_logo_audit_v2/figures/}}
\title{Integrity Treatment and Dataset Sequence Logos for GraphPro-Derived Data}
\author{CapsProm project}
\date{22 October 2026}

\begin{document}
\maketitle

\begin{abstract}
This report documents the label-agnostic integrity treatment of the released
GraphPro human and mouse collections and compares sequence logos before and
after treatment for six independently reconstructed bacterial candidates and
the two released eukaryotic integrity views. The logos describe aligned dataset
composition; they are not motifs extracted from a trained model. All training
and confirmatory-use gates remain disabled.
\end{abstract}

\section{Scope and dataset identities}

The bacterial candidates were reconstructed from Promotech TSS tables and
matching versioned assemblies. Their ``before'' version is
\texttt{graphpro-reconstructed-v1}; their ``after'' version is the
strand-invariant exact-integrity correction
\texttt{graphpro-reconstructed-v2}. Human and mouse have not yet been
independently reconstructed. Their ``before'' data are the released GraphPro
files and their ``after'' data are the exploratory integrity view
\texttt{graphpro-released-eukaryotic-integrity-v1}. This distinction prevents
computational cleaning from being misrepresented as biological label
validation.

The Eukaryotic Promoter Database (EPD) is an experimentally supported promoter
and TSS resource, not a negative-sequence catalogue
\cite{meylan2020epd}. All released human and mouse records carry
\texttt{FP}-style identifiers, including records in the directory named
\texttt{negative}. A dated EPD release and its corresponding genome assembly
must therefore be resolved before eukaryotic reconstruction. Directory names,
header suffixes, sequence composition, and agreement with another dataset were
not used to infer labels.

Table~\ref{tab:inventory} introduces the populations included in the logo
analysis. Every sequence has 81 nucleotides. The bacterial TSS coordinate is
independently positioned. Human and mouse retain the source's nominal
$-60\ldots+20$ alignment only for descriptive comparison; their TSS position
has not been independently verified.

\begin{table}[H]
\centering
\caption{Datasets and version comparisons used in this report. ``Reconstructed'' means that positives were recreated from independent TSS coordinates and negatives were newly sampled genomic controls. The eukaryotic rows are released-data integrity views, not reconstructions.}
\label{tab:inventory}
\small
\begin{tabular}{llll}
\toprule
Key & Organism & Comparison & TSS status \\
\midrule
BACILLUS & \textit{B. amyloliquefaciens} & reconstructed v1--v2 & verified \\
CJEJUNI\_2 & \textit{C. jejuni} & reconstructed v1--v2 & verified \\
CPNEUMONIAE & \textit{C. pneumoniae} & reconstructed v1--v2 & verified \\
ECOLI & \textit{E. coli} & reconstructed v1--v2 & verified \\
LINTERROGANS & \textit{L. interrogans} & reconstructed v1--v2 & verified \\
RHODOBACTER\_1 & \textit{R. capsulatus} & reconstructed v1--v2 & verified \\
human & \textit{H. sapiens} & released--integrity v1 & nominal only \\
mouse & \textit{M. musculus} & released--integrity v1 & nominal only \\
\bottomrule
\end{tabular}
\end{table}

\section{Predeclared integrity treatment}

Only criteria already established in the project were applied: (i) every
occurrence of a sequence present in both classes was removed bilaterally;
(ii) within each class, the first exact sequence in source order was retained
and subsequent occurrences were removed; and (iii) sequences not composed of
81 strict-DNA symbols were excluded. The four eukaryotic files already
contained only strict 81-nt DNA. Unresolved EPD identifiers were retained only
in this exploratory released-data view: under the reconstruction contract they
remain ineligible until resolution, but excluding all of them here would yield
an empty collection rather than a reconstructed dataset.

Table~\ref{tab:euk-treatment} reports the eukaryotic changes. Human lost one
cross-label sequence from each class, seven further positive duplicates, and
four further negative duplicates. Mouse had no cross-label collision and lost
three positive and one negative duplicate. The resulting class-count
differences are three sequences for human and two for mouse; no resampling was
used to restore balance.

\begin{table}[H]
\centering
\caption{Released human and mouse counts before and after the exploratory integrity treatment. Counts remain unbalanced after removing the observed problems because no class was downsampled.}
\label{tab:euk-treatment}
\begin{tabular}{llrrrrr}
\toprule
Organism & Class & Before & Collision & Duplicate & After & Removed \\
\midrule
Human & positive & 2,339 & 1 & 7 & 2,331 & 8 \\
Human & negative & 2,339 & 1 & 4 & 2,334 & 5 \\
Mouse & positive & 1,846 & 0 & 3 & 1,843 & 3 \\
Mouse & negative & 1,846 & 0 & 1 & 1,845 & 1 \\
\bottomrule
\end{tabular}
\end{table}

Table~\ref{tab:all-counts} presents every comparison. An unchanged row means
that the version-2 integrity filter found no affected positive--control pair;
the before and after logo matrices are consequently identical.

\begin{table}[H]
\centering
\caption{Sequence counts used in each class-specific logo before and after treatment. Bacterial removal preserves positive--control pairs; eukaryotic removal follows the released class files independently.}
\label{tab:all-counts}
\small
\begin{tabular}{lrrrrrr}
\toprule
Dataset & Pos. before & Pos. after & Removed & Neg. before & Neg. after & Removed \\
\midrule
BACILLUS & 1,064 & 1,063 & 1 & 1,064 & 1,063 & 1 \\
CJEJUNI\_2 & 1,905 & 1,882 & 23 & 1,905 & 1,882 & 23 \\
CPNEUMONIAE & 529 & 529 & 0 & 529 & 529 & 0 \\
ECOLI & 248 & 248 & 0 & 248 & 248 & 0 \\
LINTERROGANS & 2,783 & 2,783 & 0 & 2,783 & 2,783 & 0 \\
RHODOBACTER\_1 & 5,374 & 5,321 & 53 & 5,374 & 5,321 & 53 \\
Human & 2,339 & 2,331 & 8 & 2,339 & 2,334 & 5 \\
Mouse & 1,846 & 1,843 & 3 & 1,846 & 1,845 & 1 \\
\bottomrule
\end{tabular}
\end{table}

\section{Sequence-logo method}

The figures use the sequence-logo construction of Schneider and Stephens
\cite{schneider1990logos}. At each aligned position, letter heights equal the
observed nucleotide frequency multiplied by $2-H$, where $H$ is the base-2
Shannon entropy of the four nucleotide frequencies. No pseudocount or
small-sample correction was applied. Every panel uses positions $-60$ through
$+20$, a fixed 0--2 bit scale, identical nucleotide colours, and the same
15-by-6 inch rendering. The dashed line marks the verified TSS for bacteria and
the explicitly labelled nominal TSS for human and mouse. Frequency and
information-height matrices are stored as CSV files alongside the figures.

Figures~\ref{fig:bacillus}--\ref{fig:mouse} show the class-specific dataset
logos. Each embedded image contains its before panel above its after panel and
states the number of contributing sequences. These are descriptive dataset
logos and must not be cited as motifs learned by CapsProm or any baseline.

\newcommand{\datasetlogos}[5]{%
\begin{figure}[H]
\centering
\includegraphics[width=0.96\textwidth]{#1_positive_before_after.png}
\includegraphics[width=0.96\textwidth]{#1_negative_before_after.png}
\caption{#2 dataset sequence logos before and after #3. The positive class is shown first and the negative class second. All panels use the same coordinate and information scales. #4}
\label{#5}
\end{figure}}

\datasetlogos{bacillus}{\textit{Bacillus amyloliquefaciens}}{pair-preserving exact-integrity filtering}{One repeated positive and its paired control were removed.}{fig:bacillus}
\datasetlogos{cjejuni_2}{\textit{Campylobacter jejuni}}{pair-preserving exact-integrity filtering}{Twenty-three repeated positives and their paired controls were removed.}{fig:cjejuni}
\datasetlogos{cpneumoniae}{\textit{Chlamydophila pneumoniae}}{pair-preserving exact-integrity filtering}{No sequence was removed; before and after matrices are identical.}{fig:cpneumoniae}
\datasetlogos{ecoli}{\textit{Escherichia coli}}{pair-preserving exact-integrity filtering}{No sequence was removed; before and after matrices are identical.}{fig:ecoli}
\datasetlogos{linterrogans}{\textit{Leptospira interrogans}}{pair-preserving exact-integrity filtering}{No sequence was removed; before and after matrices are identical.}{fig:linterrogans}
\datasetlogos{rhodobacter_1}{\textit{Rhodobacter capsulatus}}{pair-preserving exact-integrity filtering}{Fifty-three repeated positives and their paired controls were removed.}{fig:rhodobacter}
\datasetlogos{human}{\textit{Homo sapiens}}{released-data integrity filtering}{The coordinate at zero is nominal and not independently verified.}{fig:human}
\datasetlogos{mouse}{\textit{Mus musculus}}{released-data integrity filtering}{The coordinate at zero is nominal and not independently verified.}{fig:mouse}

\section{Comparative observations}

The removals have negligible visual impact. Across all changed populations,
the largest absolute shift in aggregate GC fraction is approximately 0.001 and
the largest shift in mean positional information is below 0.001 bits. This is
consistent with removing at most 53 of 5,374 sequences in any bacterial class
and at most eight of 2,339 sequences in an eukaryotic class. Thus, the treatment
removes integrity violations without materially redefining the gross
composition of the retained collections.

The bacterial logos generally show low per-position information. CJEJUNI\_2
has the largest mean positional information among these populations, whereas
ECOLI is nearly compositionally diffuse on this aggregate scale. The paired
genomic controls were GC-matched within four nucleotides per pair, but the logos
do not establish that every control is biologically inactive. Nor does a weak
aggregate logo imply that a classifier cannot exploit distributed or
nonlinear signals.

The released human and mouse classes display a substantial class-level GC
difference: approximately 57.3\% versus 46.4\% for human and 56.1\% versus
45.7\% for mouse. This is a potential dataset-construction shortcut and merits
formal distribution testing after provenance resolution. A visibly conserved
terminal position is also present in the released eukaryotic alignments. Since
their assembly, coordinate extraction, and nominal TSS have not been
independently verified, this terminal pattern is reported as a source-alignment
observation, not as evidence of a promoter element.

Canonical promoter elements such as TATA-like patterns or initiator-associated
composition may motivate later, predeclared enrichment tests. Visual inspection
alone cannot establish motif significance, functional activity, or causal
importance. Such claims require background-aware statistics, multiple-testing
control, independent annotation, and preferably replication across sources.

\section{Reproducibility and remaining gates}

The eukaryotic treatment contract is
\texttt{configs/graphpro\_eukaryotic\_integrity\_v1.yaml}; its builder writes a
row-level exclusion ledger with dataset, organism, released class, source
index, identifier, description, sequence hash, and exclusion reason. The final
logo contract is
\texttt{configs/graphpro\_dataset\_logo\_audit\_v2.yaml}. The v1 logo output is
preserved because it labelled the eukaryotic nominal coordinate too strongly;
v2 corrects only that presentation and does not change matrices or populations.

From the repository root, the products can be recreated with:
\begin{verbatim}
PYTHONPATH=src .venv/bin/python \
  scripts/build_graphpro_eukaryotic_integrity.py \
  --output-dir data/interim/graphpro_eukaryotic_integrity_v1
MPLCONFIGDIR=/tmp/capsprom-mpl PYTHONPATH=src .venv/bin/python \
  scripts/generate_graphpro_dataset_logos.py \
  --config configs/graphpro_dataset_logo_audit_v2.yaml \
  --output-dir outputs/reports/graphpro_dataset_logo_audit_v2
\end{verbatim}

Human and mouse reconstruction remains blocked until a dated EPD release,
assembly, identifier-resolution table, coordinates, and independently generated
controls are frozen. For all organisms, remaining work includes primary-source
sampling, distribution and cross-dataset dependency audits, licensing metadata,
and independent manual review. Existing model results tied to released GraphPro
or canonical-v3 identities remain historical exploratory provenance and cannot
be reused as results for these new identities. No training, final-fit,
external-test, embedding, or confirmatory gate was opened for this report.

\begin{thebibliography}{9}
\bibitem{schneider1990logos}
T. D. Schneider and R. M. Stephens.
Sequence logos: a new way to display consensus sequences.
\emph{Nucleic Acids Research}, 18(20):6097--6100, 1990.
doi:10.1093/nar/18.20.6097.

\bibitem{meylan2020epd}
P. Meylan et al.
EPD in 2020: enhanced data visualization and extension to ncRNA promoters.
\emph{Nucleic Acids Research}, 48(D1):D65--D69, 2020.
doi:10.1093/nar/gkz1014.
\end{thebibliography}

\end{document}
